\scanpage{175}
il faut et il suffit que l'on ait $\widetilde C\omega=\widetilde\omega\otimes_{\underline O_{X^{p/S}}}1_{\underline\Lambda}$.

La première assertion est \unsure{standart|standard}, et ne dépend pas de l'hyp. de car. $p>0$. [Si $\underline L$ n'est pas de rang 1, mais seulement loc. libre, alors $\omega$ est une forme diff. à valeurs dans $\underline{\operatorname{End}}(E)$, et la relation à écrire est $d\omega+[\omega,\omega]=0$]

% slanted
\margin{\uline{maladroit de choisir des bases} Dire : la section de $(\Omega^1_{X/S})^{\otimes p}$ ``défaut de compatibilité aux puiss. $p^{\text{ièmes}}$'' de la connexion envisagée $\Delta$ [ou $p$-courbure] est égale à [celle de $\Delta_0$, supposée nulle, plus] $\omega^{\otimes p}-f_{X/S}^*\widetilde C\omega=f_{X/S}^*(\widetilde\omega-\widetilde C\omega)$ donc est nulle ssi on a $\widetilde\omega=\widetilde C\omega$ \ldots.}

Pour la deuxième, la condition étant locale, on est ramené au cas où tout est affine. On a donc, si on choisit une $p$-base $(x_i)$, d'où des $D_i=\frac{\partial}{\partial x_i}$,
\deleted{$\unsure{D_\xi,\rho_\Delta}(\unsure{D})\quad\rho(D_i)=\rho_0(D_i)+\unsure{u}$}
\deleted{d'où} Ecrivons la condition que $\Delta$, \deleted{\unsure{associé à}}\inserted{qui donne} $\xi\longrightarrow\Theta_\xi$ \deleted{$\Theta(D_i)^p$} est compatible aux puissances $p$-ièmes. On trouve \deleted{$\unsure{\Theta_{D_i}^p}$} les conditions nécessaires
\[
 \Theta_{D_i}^{p}=0\qquad i\in \unsure{I|J}
\]
car $D_i^p=0$. Je dis qu'elle est suffisante. \deleted{\unsure{En}}\inserted{Cela} se voit aisément, en prouvant que l'algèbre des \deleted{\unsure{effet, tte dérivation est de la forme $\sum a_iD_i$,}} op. différentiels engendrés par les $D_i=\frac{\partial}{\partial x_i}$ et les mult. \deleted{les $a_i\in B$.} scalaires est engendrée par \inserted{les générateurs $D_i,x_i$, avec} les relations
\margin{Indépendant de l'hyp. que $L$ de rang 1}
\[
 \left\{\begin{aligned}
 [D_i,D_j]&=0\\
 [x_i,x_j]&=0\\
 [D_i,x_j]&=\delta_{ij}\\
 x_i^p&=a_i\\
 D_i^p&=0
 \end{aligned}\right.
\]
% slanted
\margin{\uline{mieux} : définir l'obstruction à la compatibilité aux puiss. $p$.ièmes comme une section de $(\Omega^1_{X/S})^{\otimes p}\otimes\underline{\operatorname{End}}(E)$ (pour une connexion \uline{à courbure nulle} donnée)}

Il faut donc, compte tenu de $\overset{\Delta_i}{\Theta_{D_i}}=\overset{\Delta_i^0}{\Theta^0_{D_i}}+\omega(D_i).\mathrm{id}$ écrivant $\omega=\sum c_i\,dx_i$, où $c_i=\omega(D_i)$, exprimer la condition
\[
 (\Delta_i^0+c_i\cdot)^p=0
\]
Mais la formule de Jacobson donne
% archive sheet number: 174

